\chapter{METHODOLOGY}
\label{ch:method}

\section{Maintenance Pipeline}

A maintenance run turns one development push into a reviewable knowledge change in six stages.
\begin{enumerate}[nosep]
  \item \textbf{Collection.} The runtime freezes the inputs: base and head commits of the application, the knowledge baseline commit, the selected Claude Code sessions, and a timezone-aware cutoff. Selected conversations are normalized (roles, order, and message identifiers are kept; tool results and model thinking are dropped) and hashed, so later runs can verify that they read the same evidence.
  \item \textbf{Evidence pack.} The explicit diff, the list of changed files, the normalized conversations, and optional memory files are written to a read-only directory that every agent of the run reads.
  \item \textbf{Initialization.} Plain code reads the committed configuration (\code{init.md}) and selects the specialists. The four core specialists always run. A tool specialist runs only when the push touches files matching its folder's \code{fileMatchPattern} or mentions its keywords.
  \item \textbf{Investigation and proposals.} Specialists investigate the evidence and return structured change records (Section~\ref{sec:record}). Strategies A and B differ only in this stage (Section~\ref{sec:strategies}).
  \item \textbf{Handoffs and review.} A specialist that received handoffs from others gets one more turn on a fork of its own session to act on them. Deterministic checks then validate every record: its target must belong to the specialist, its operation must be valid, and its evidence locators must resolve to a real session message, diff path, commit log, or source line. Unresolvable locators are dropped; a record with none left is rejected, and records whose support is unclear are deferred. An optional cartographer proposes new documents for supported findings that no existing document can hold.
  \item \textbf{Writing.} One writer agent applies the accepted records to a temporary worktree of the knowledge repository, appends a log entry, and commits to a new local branch. The runtime reports the actual Git diff, independently of what the writer claims it applied.
\end{enumerate}

\section{Change Records and Evidence}
\label{sec:record}

Every proposal is a JSON record with a target document, a fact identifier, an operation (ADD, UPDATE, DELETE, MERGE, NO\_CHANGE, or NEEDS\_REVIEW), the final wording, a scope, a reason, and at least one evidence locator. Locators point to a conversation message (session and message identifier with a short verbatim quote), a diff path, a source line at the pushed head, or a memory file. The evidence rules separate kinds of support: source code establishes what is implemented; developer statements establish intent, definitions, requirements, and decisions; an agent's completion message alone establishes nothing. Later messages override earlier ones, so a specialist must read a session to its end before admitting a fact. A choice observed in one application is scoped to that application and never promoted to a company standard without an explicit statement.

\section{Context-Acquisition Strategies}
\label{sec:strategies}

\textbf{Strategy A (sweeper first)} starts one sweeper agent that reads all of the evidence and records candidate findings for each specialist. When the sweeper finishes, the runtime forks its completed session once per specialist using the SDK's session-resume-and-fork mechanism. Each fork inherits the sweeper's entire context and then receives its own role. If the sweeper fails, the specialists do not run; A never falls back to B.

\textbf{Strategy B (fresh specialists)} starts every specialist in a new session with only its role and the evidence pack. Each specialist decides for itself what to read.

Both strategies use the same role definitions, rules, tools, evidence, review, and writer, so a difference in outcome can be attributed to how context is acquired.

\paragraph{Prompt-cache sharing.} A fork is only cheap if the model provider can reuse the cached prefix of the sweeper's conversation. The cache prefix is ordered tool definitions, then system prompt, then messages, and any difference invalidates everything after it. In the pilot every fork received a different system prompt and output schema from its sweeper, so each fork re-wrote about 42{,}000 inherited tokens to the cache, and A cost 2.2 times as much as B. We changed the runtime so that the sweeper and every specialist share one role-free system prompt and one output schema; the role arrives as the first user message after the inherited history. A live run then showed a second cause: connectors attached to the user's login connected asynchronously and changed the tool list between calls. Maintenance agents now run without MCP servers. After both changes every fork call read the inherited history from the cache.

\section{Simulation of Development}
\label{sec:simulation}

Evaluating maintenance requires development conversations with known content. Real conversations are private and have no ground truth, so we generate them. For each application we wrote task briefs that contain the task, the expected result, and \emph{moments}: facts a developer would state during the work, such as a correction (``store amounts at the same scale for every currency''), a decision (``use undici's fetch on the server''), or a definition. Each brief also contains decoys, statements that look like knowledge but must not be admitted (for example a temporary workaround), and deviations.

A persona model plays the developer and talks to a real Claude Code agent working in a real checkout of the application. In v2 each application has three teammate personas with their own seniority and voice. Personas open with vague requests, iterate on the agent's work, ask questions, and must deliver every pending moment in their own words before they allow a commit; a filter rejects messages that mention knowledge, documentation, evaluation, or other words that would reveal the staging. The commits and sessions are then frozen as the input of every maintenance condition.

\section{Evaluation of Maintenance (RQ1)}

For each push, a blind reviewer model receives the expected facts, the decoys, and the accepted records of each condition under shuffled labels. It marks each required fact as found, partial, or missed, and each decoy as admitted or not. Facts that the persona never actually delivered are excluded. We report required facts found, partial matches, optional facts, traps admitted, estimated cost, and wall-clock time per run. The writer's fidelity to accepted records was not scored in v2; proposals and applied documents are separate measurements.

\section{Evaluation of Downstream Use (RQ2)}
\label{sec:downstream-method}

The downstream study asks whether the maintained knowledge helps a development agent. The protocol was written and frozen before any measured run.

\paragraph{Conditions.} All conditions start from the same application revision (the head that the v2 maintenance runs saw).
\begin{itemize}[nosep]
  \item \textbf{none}: no knowledge and no pointer to it; the agent sees the source and its Git history.
  \item \textbf{seed}: \code{CLAUDE.md} points to the v2 baseline knowledge, before maintenance.
  \item \textbf{maintained (A)}: the same pointer to the baseline plus the knowledge patch written by the v2 A-Sonnet maintenance run.
  \item \textbf{maintained (B)}: the baseline plus the patch of the v2 B-Sonnet run. This condition was added after the first three had been analysed, to test whether A's better maintenance carries over to development.
\end{itemize}
The knowledge directories differ only by those patches. They are read-only and have opaque names, so the agent cannot tell which revision it reads.

\paragraph{Tasks.} Each application has two tasks (Appendix~\ref{app:tasks}). A \emph{routine} task can be finished from the code alone; a \emph{knowledge} task has review checks that lean on project decisions. Prompts are short, informal teammate messages that never name the checked conventions. Each task has 6 to 8 hidden review checks written before execution. Each check has a \emph{basis}, used only in the analysis: \emph{code} (visible in the source or history), \emph{seed} (stated in the baseline knowledge), \emph{maintained} (added or corrected by the maintenance patch), or \emph{unstated} (true in the project but recorded in no knowledge revision).

\paragraph{Execution.} The development agent is Claude Code (Sonnet 5.5) with its standard system prompt, project settings only, file and shell tools, no web access, no subagents, and no dependency installation. Each run gets a fresh Git worktree. If the agent stops without committing, it receives one fixed neutral reply, at most twice; the number of extra messages is the clarification count. Each (task, condition) cell is repeated five times in an interleaved fixed order: 90 runs for the first three conditions and 30 more for maintained (B).

\paragraph{Scoring.} A blind judge (Opus 5.5) sees only the task, the checklist, the commit messages, and the diff. It marks each check pass, fail, or not applicable, rates completion from 0 to 3, and lists defects. The primary quality metric is the pass rate over applicable checks. We also record cost, wall-clock time, turns, tool calls, knowledge files read, and clarification rounds.

\section{Evaluation over a Feature Sequence (RQ2, end to end)}
\label{sec:loop-method}

The downstream comparison gives the agent knowledge produced earlier from simulated conversations. The feature-sequence experiment instead runs the complete loop of Figure~\ref{fig:loop} and asks whether decisions a developer states while working reach later features developed in fresh sessions. The protocol was frozen before the measured runs.

\paragraph{Features and decisions.} Each application has a chain of three small features in one area of the product, developed in order on the same checkout (Appendix~\ref{app:loop}). F1's prompt states two team decisions in passing (``heads up, we agreed that from now on \dots''), F2's prompt states two more, and F3's prompt states none. Linkwarden's F1 also states a commit-message convention (Conventional Commits). Every decision was chosen so that the existing code does not settle it: the code has mixed patterns, or the decision concerns something the code shape does not show, such as where previews go or how errors leak information. Each feature has 7 to 9 hidden checks with a basis: \emph{code}, \emph{stated} (a decision in the same prompt), \emph{carried} (a decision from an earlier prompt), or, in Actual Budget and Maybe, \emph{leak} (Linkwarden's commit convention must not appear).

\paragraph{Conditions.}
\begin{itemize}[nosep]
  \item \textbf{none}: no knowledge. Later features see the code and history left by earlier ones, but no earlier conversation.
  \item \textbf{repo}: one empty knowledge repository per application (\code{scope: repository}) on a local bare remote, connected with \code{contextops attach}. After F1 and F2, the real pre-push hook runs strategy A on the pushed commits and the sessions since the last run and pushes its result branch; a simulated reviewer merges every branch unchanged.
  \item \textbf{team}: one empty knowledge repository shared by the three applications (\code{scope: team}). Each round, the three applications develop the next feature against the same merged knowledge, then push in a fixed order.
\end{itemize}
Each condition is repeated three times: 81 features, 36 maintenance runs. The development agent, tools, limits, clarification policy, and blind judge are those of the downstream comparison; the agent loads project and local settings so that the attach pointer applies. The primary outcome is the pass rate of carried checks in F2 and F3.

\section{Maintenance Rerun and Persona-Driven Development (v3)}
\label{sec:v3-method}

\paragraph{Maintenance rerun.} We reran the four v2 conditions (A or B, Sonnet or Haiku) on exactly the v2 inputs: the same push and the same six conversations per application, the same cutoff, the same seed knowledge, and v2's budgets. Only the runtime changed: the fixes found by the live first runs, the application scope, \code{scope: repository}, and the admission policy. One blind judge (Opus 5.5) scored the accepted records of the v3 and the v2 runs together, as eight anonymous labels per application, against the v2 expected facts and traps.

\paragraph{Development with personas.} Each v3 result (seed plus result branch) became a knowledge base for a development agent connected with \code{contextops attach}; a fifth condition had no knowledge. Agents designed two new tasks per application (six in total) that touch two to five facts from the v2 conversations, preferring facts on which the v2 conditions differed. A simulated teammate (Sonnet 5.5, a named v2 persona) opens with a short natural request, answers the agent's questions from its private knowledge, and after a commit skims the diff and may ask for fixes on its review focus, at most twice. The judge scores the diff at the agent's first commit (the agent alone) and the final diff (after review); we also count review rounds and the facts the teammate had to tell. Three repetitions give 90 sessions. After the first analysis we added two conditions with a session-start digest of the knowledge (Section~\ref{sec:v3-dev}), 36 more sessions; they are reported as a post hoc amendment.

\section{Team Decisions Across Repositories}
\label{sec:v4-method}

The last experiment tests the use the repository is built for: a team decision stated in one repository and needed in another, where the code shows nothing about it. The protocol was written before any session.

\paragraph{Learning.} One empty team knowledge repository (\code{scope: team}) on a local remote is attached to a Linkwarden checkout with the pre-push hook and the session-start digest. In three persona sessions, a tech lead asks for small Linkwarden changes and states four decisions in passing as rules for all of the team's repositories: no new third-party package for small things (TD1); no personal data or secrets in logs at any level (TD2); every outbound call has an explicit timeout of at most 10 seconds and its failure never breaks the user's action (TD3); commit subjects follow Conventional Commits (TD4). Each push runs maintenance and the result is merged unchanged. A design agent checked in the code that Actual Budget and Maybe do not reveal these rules; Actual Budget's recent history uses a different commit prefix and Maybe's code logs user emails and sets no HTTP timeouts.

\paragraph{Testing.} Five new tasks, two in Actual Budget, two in Maybe and one in Linkwarden, each give a natural chance to break two or three of the decisions (an \code{.ics} export that tempts a calendar package, a webhook on account creation, a CSV encoding fix, an exchange-rate fallback, short-link expansion). Persona teammates open naturally and never state the decisions; they know them privately and may ask for fixes in review. Conditions: no knowledge, or the frozen team knowledge attached with the digest. Four repetitions give 40 sessions. The primary outcome is the pass rate of decision checks at the agent's first commit.

\section{Reproducible First Run}

To show the template-to-populated path to outside users, we prepared public inputs for one application: the upstream Actual Budget repository at a pinned commit, a Git bundle with the eight development commits of the v2 push (re-authored with a neutral identity and with local paths removed), and the six development sessions exported with the new \code{contextops export-session} command. We then followed the documented first run from a clean clone: install, run the tests, initialize empty knowledge with \code{contextops init}, import the sessions, run strategy A, and inspect the result.
